\textbf{1. Kiến thức cần nhớ}
\begin{itemize}
\item \textbf{Định nghĩa}: Tứ giác nội tiếp là tứ giác có bốn đỉnh cùng nằm trên một đường tròn.
\item \textbf{Định lí}: Trong một tứ giác nội tiếp, \textbf{tổng hai góc đối diện bằng $180^\circ$}:
$$\boxed{\widehat{A} + \widehat{C} = 180^\circ;\quad \widehat{B} + \widehat{D} = 180^\circ}$$
\item \textbf{Định lí đảo}: Nếu một tứ giác có tổng hai góc đối diện bằng $180^\circ$ thì tứ giác đó nội tiếp được đường tròn.
\item \textbf{BỐN DẤU HIỆU NHẬN BIẾT TỨ GIÁC NỘI TIẾP (CỰC KÌ QUAN TRỌNG CHO THI TUYỂN SINH 10)}:
\begin{enumerate}
\item \textbf{Dấu hiệu 1 (Tổng hai góc đối bằng $180^\circ$)}: Tứ giác có $\widehat{A} + \widehat{C} = 180^\circ$ (đặc biệt: hai góc vuông đối diện $\widehat{B} + \widehat{D} = 90^\circ + 90^\circ = 180^\circ$).
\item \textbf{Dấu hiệu 2 (Góc ngoài tại một đỉnh bằng góc trong tại đỉnh đối)}: $\widehat{A} = \widehat{C}_{\text{ngoài}}$.
\item \textbf{Dấu hiệu 3 (Cùng nhìn một cạnh dưới góc bằng nhau)}: Hai đỉnh kề nhau cùng nhìn cạnh chứa hai đỉnh còn lại dưới một góc bằng nhau ($\widehat{DAC} = \widehat{DBC}$).
\item \textbf{Dấu hiệu 4 (Cách đều một điểm)}: Bốn đỉnh cùng cách đều một điểm $O$ cố định ($OA = OB = OC = OD$).
\end{enumerate}
\end{itemize}

\textbf{2. Phương pháp giải các dạng}

%% DANG 9H9?29-1
\textbf{Dạng 1. Chứng minh tứ giác nội tiếp}
\begin{enumerate}
\item Nhìn hình nhận diện dấu hiệu:
\begin{itemize}
\item Nếu có hai góc vuông đối diện $\rightarrow$ dùng Dấu hiệu 1.
\item Nếu có hai góc vuông cùng nhìn một cạnh $\rightarrow$ dùng Dấu hiệu 3 (đường kính là cạnh chung đó).
\item Nếu có các đường chéo cắt nhau tạo thành các góc bằng nhau $\rightarrow$ dùng Dấu hiệu 3.
\item Nếu có tia đối của một cạnh $\rightarrow$ dùng Dấu hiệu 2 (góc ngoài bằng góc trong đối).
\end{itemize}
\item Trình bày lập luận chặt chẽ, chỉ rõ tên góc và lí do trước khi kết luận tứ giác nội tiếp.
\end{enumerate}

%% DANG 9H9?29-2
\textbf{Dạng 2. Tính góc trong tứ giác nội tiếp}
\begin{enumerate}
\item Sử dụng tính chất: góc đối bù nhau $\widehat{C} = 180^\circ - \widehat{A}$.
\item Chuyển góc: hai góc nội tiếp cùng chắn một cung trong đường tròn ngoại tiếp tứ giác thì bằng nhau ($\widehat{ADB} = \widehat{ACB}$).
\end{enumerate}

%% DANG 9H9?29-3
\textbf{Dạng 3. Bài toán tổng hợp về tứ giác nội tiếp (Câu hình thi Tuyển sinh 10)}
\begin{enumerate}
\item \textbf{Chứng minh song song / vuông góc}: Dùng tứ giác nội tiếp để chuyển góc $\rightarrow$ suy ra hai góc so le trong/đồng vị bằng nhau.
\item \textbf{Chứng minh các hệ thức tích độ dài}: Từ tứ giác nội tiếp suy ra hai tam giác đồng dạng $\rightarrow$ lập tỉ số đồng dạng suy ra hệ thức tích (ví dụ: $\triangle MAB \sim \triangle MCD \implies MA \cdot MB = MC \cdot MD$).
\item \textbf{Chứng minh ba điểm thẳng hàng, các đường thẳng đồng quy}: Vận dụng trực tâm tam giác, điểm chung của các đường tròn ngoại tiếp hoặc tính chất góc bẹt ($180^\circ$).
\end{enumerate}
